\chapter{2013年湖北卷（理）}

\section{单选题}

\begin{problem}
在复平面内，复数 \(z = \frac{2\i}{1 + \i}\)（\(\i\) 为虚数单位）的共轭复数对应的点位于（\quad）

\choices
    {第一象限}
    {第二象限}
    {第三象限}
    {第四象限}
\end{problem}

\begin{answer}
D
\end{answer}

\begin{solution}
将分子分母同乘以 \(1-\i\)，得
\(z=\dfrac{2\i(1-\i)}{(1+\i)(1-\i)}=\dfrac{2\i(1-\i)}{2}=1+\i\).
因此共轭复数为 \(\overline{z}=1-\i\)，它在复平面内对应的点为 \((1,-1)\)，位于第四象限，故选 D.
\end{solution}

\begin{problem}
已知全集为 \(\R\). 集合 \(A = \left\{x \mid \left(\frac{1}{2}\right)^x \leqslant 1\right\}\)，\(B = \{x \mid x^2 - 6x + 8 \leqslant 0\}\)，则 \(A \cap \complement_{\R} B =\)（\quad）

\choices
    {\(\{x \mid x \leqslant 0\}\)}
    {\(\{x \mid 2 \leqslant x \leqslant 4\}\)}
    {\(\{x \mid 0 \leqslant x < 2\}\) 或 \(\{x \mid x > 4\}\)}
    {\(\{x \mid 0 < x \leqslant 2\}\) 或 \(\{x \mid x \geqslant 4\}\)}
\end{problem}

\begin{answer}
C
\end{answer}

\begin{solution}
因为函数 \(y=\left(\dfrac12\right)^x\) 是减函数，且在 \(x=0\) 时取值为 \(1\)，所以 \(\left(\dfrac12\right)^x\leqslant1\) 等价于 \(x\geqslant0\)，即 \(A=[0,+\infty)\).
又 \(x^2-6x+8=(x-2)(x-4)\)，所以 \(B=[2,4]\)，从而 \(\complement_{\R}B=(-\infty,2)\cup(4,+\infty)\).
因此
\(A\cap\complement_{\R}B=[0,2)\cup(4,+\infty)\)，对应选项 C.
\end{solution}

\begin{problem}
在一次跳伞训练中，甲，乙两位学员各跳一次. 设命题 \(p\) 是“甲降落在指定范围”，\(q\) 是“乙降落在指定范围”，则命题“至少有一位学员没有降落在指定范围”可表示为（\quad）

\choices
    {\((\neg p) \vee (\neg q)\)}
    {\(p \vee (\neg q)\)}
    {\((\neg p) \wedge (\neg q)\)}
    {\(p \vee q\)}
\end{problem}

\begin{answer}
A
\end{answer}

\begin{solution}
“至少有一位学员没有降落在指定范围”是“不是两位学员都降落在指定范围”，即命题 \(\neg(p\wedge q)\).
由德摩根律，\(\neg(p\wedge q)=(\neg p)\vee(\neg q)\)，所以选 A.
\end{solution}

\begin{problem}
将函数 \( y = \sqrt{3} \cos x + \sin x \)（\(x \in \R\)）的图像向左平移 \(m\)（\(m > 0\)）个单位长度后，所得到的图象关于 \(y\) 轴对称. 则 \(m\) 的最小值是（\quad）

\choices
    {\(\frac{\pi}{12}\)}
    {\(\frac{\pi}{6}\)}
    {\(\frac{\pi}{3}\)}
    {\(\frac{5\pi}{6}\)}
\end{problem}

\begin{answer}
B
\end{answer}

\begin{solution}
原函数可化为
\(y=\sqrt{3}\cos x+\sin x=2\cos\left(x-\dfrac{\pi}{6}\right)\).
图象向左平移 \(m\) 个单位后，函数变为
\(g(x)=2\cos\left(x+m-\dfrac{\pi}{6}\right)\).
图象关于 \(y\) 轴对称等价于 \(g(x)\) 为偶函数. 展开得
\(g(x)=2\cos x\cos\left(m-\dfrac{\pi}{6}\right)-2\sin x\sin\left(m-\dfrac{\pi}{6}\right)\)，故必须有
\(\sin\left(m-\dfrac{\pi}{6}\right)=0\).
于是 \(m=\dfrac{\pi}{6}+k\pi\)，其中 \(k\) 为整数. 结合 \(m>0\)，最小值为 \(\dfrac{\pi}{6}\)，故选 B.
\end{solution}

\begin{problem}
已知 \(0 < \theta < \frac{\pi}{4}\)，则双曲线 \(C_1: \frac{x^{2}}{\cos^{2}\theta} - \frac{y^{2}}{\sin^{2}\theta} = 1\) 与 \(C_2: \frac{y^{2}}{\sin^{2}\theta} - \frac{x^{2}}{\sin^{2}\theta \tan^{2}\theta} = 1\) 的（\quad）

\choices
    {实轴长相等}
    {虚轴长相等}
    {焦距相等}
    {离心率相等}
\end{problem}

\begin{answer}
D
\end{answer}

\begin{solution}
双曲线 \(C_1\) 的实半轴和虚半轴分别为 \(a_1=\cos\theta\)、\(b_1=\sin\theta\)，所以其焦半距满足 \(c_1^2=a_1^2+b_1^2=1\)，离心率为
\(e_1=\dfrac{c_1}{a_1}=\dfrac{1}{\cos\theta}\).
将 \(C_2\) 写成标准形式
\(\dfrac{y^2}{(\sin\theta)^2}-\dfrac{x^2}{(\sin\theta\tan\theta)^2}=1\)，其竖直实半轴为 \(a_2=\sin\theta\)，虚半轴为 \(b_2=\sin\theta\tan\theta\).
因此
\(c_2^2=a_2^2+b_2^2=\sin^2\theta+\sin^2\theta\tan^2\theta=\tan^2\theta\)，从而 \(c_2=\tan\theta\)，且
\(e_2=\dfrac{c_2}{a_2}=\dfrac{\tan\theta}{\sin\theta}=\dfrac{1}{\cos\theta}=e_1\).
故两条双曲线的离心率相等，选 D.
\end{solution}

\begin{problem}
已知点 \(A(-1,1)\)，\(B(1,2)\)，\(C(-2,-1)\)，\(D(3,4)\). 则向量 \(\overrightarrow{AB}\) 在 \(\overrightarrow{CD}\) 方向上的投影为（\quad）

\choices
    {\(\frac{3\sqrt{2}}{2}\)}
    {\(\frac{3\sqrt{15}}{2}\)}
    {\(-\frac{3\sqrt{2}}{2}\)}
    {\(-\frac{3\sqrt{15}}{2}\)}
\end{problem}

\begin{answer}
A
\end{answer}

\begin{solution}
\(\overrightarrow{AB}=(1-(-1),2-1)=(2,1)\)，\(\overrightarrow{CD}=(3-(-2),4-(-1))=(5,5)\).
向量 \(\overrightarrow{AB}\) 在 \(\overrightarrow{CD}\) 方向上的投影为
\(\dfrac{\overrightarrow{AB}\cdot\overrightarrow{CD}}{|\overrightarrow{CD}|}=\dfrac{2\cdot5+1\cdot5}{\sqrt{5^2+5^2}}=\dfrac{15}{5\sqrt{2}}=\dfrac{3\sqrt{2}}{2}\)，故选 A.
\end{solution}

\begin{problem}
一辆汽车在高速公路上行驶，由于遇到紧急情况而刹车，以速度 \(v(t) = 7 - 3t + \frac{25}{1 + t}\)（\(t\) 的单位：\(\mathrm{s}\)，\(v\) 的单位：\(\mathrm{m/s}\)）行驶至停止，在此期间汽车继续行驶的距离（单位：\(\mathrm{m}\)）是（\quad）

\choices
    {\(1 + 25\ln 5\)}
    {\(8 + 25\ln \frac{11}{3}\)}
    {\(4 + 25\ln 5\)}
    {\(4 + 50\ln 2\)}
\end{problem}

\begin{answer}
C
\end{answer}

\begin{solution}
汽车停止时速度为 \(0\). 由于 \(t\geqslant0\)，由
\(7-3t+\dfrac{25}{1+t}=0\) 得 \((7-3t)(1+t)+25=0\)，即 \(3t^2-4t-32=0\).
解得 \(t=4\) 或 \(t=-\dfrac{8}{3}\)，故汽车在 \(0\leqslant t\leqslant4\) 内行驶.
行驶距离为
\[
\int_0^4\left(7-3t+\frac{25}{1+t}\right)\,\mathrm{d}t
=\left[7t-\frac{3}{2}t^2+25\ln(1+t)\right]_0^4
=4+25\ln5
\]
故选 C.
\end{solution}

\begin{problem}
一个几何体的三视图如图所示，该几何体从上到下由四个简单几何体组成，其体积分别记为 \(V_{1}\)，\(V_{2}\)，\(V_{3}\)，\(V_{4}\). 上面两个简单几何体均为旋转体，下面两个简单几何体均为多面体，则有（\quad）

\bitmapfigure[float=inline,max width=0.46\linewidth]{img/2013/hubei_science/q8_fig1.png}
\FloatBarrier

\choices
    {\(V_{1} < V_{2} < V_{4} < V_{3}\)}
    {\(V_{1} < V_{3} < V_{2} < V_{4}\)}
    {\(V_{2} < V_{1} < V_{3} < V_{4}\)}
    {\(V_{2} < V_{3} < V_{1} < V_{4}\)}
\end{problem}

\begin{answer}
C
\end{answer}

\begin{solution}
由三视图可读出，从上到下四个几何体分别是高为 \(1\)、上、下底面半径分别为 \(2\)、\(1\) 的圆台，高为 \(2\)、半径为 \(1\) 的圆柱，边长为 \(2\) 的正方体，以及高为 \(1\)、上、下底面边长分别为 \(2\)、\(4\) 的正四棱台.
因此
\(V_1=\dfrac{\pi}{3}(2^2+2\cdot1+1^2)\cdot1=\dfrac{7\pi}{3}\)，
\(V_2=\pi\cdot1^2\cdot2=2\pi\)，
\(V_3=2^3=8\)，
\(V_4=\dfrac{1}{3}(2^2+2\cdot4+4^2)\cdot1=\dfrac{28}{3}\).
利用 \(\pi<\dfrac{22}{7}\)，有 \(V_2<V_1<8=V_3<V_4\)，故选 C.
\end{solution}

\begin{problem}
如图，将一个各面都涂了油漆的正方体切割为 125 个同样大小的小正方体，经过搅拌后，从中随机取一个小正方体，记其涂漆面数为 \(X\)，则 \(X\) 的均值 \(E(X) =\)（\quad）

\bitmapfigure[float=inline,max width=0.42\linewidth]{img/2013/hubei_science/q9_fig1.png}

\choices
    {\(\frac{126}{125}\)}
    {\(\frac{6}{5}\)}
    {\(\frac{168}{125}\)}
    {\(\frac{7}{5}\)}
\end{problem}

\begin{answer}
B
\end{answer}

\begin{solution}
正方体被切成 \(5^3=125\) 个小正方体. 三面涂漆的小正方体有 \(8\) 个，两面涂漆的有 \(12(5-2)=36\) 个，一面涂漆的有 \(6(5-2)^2=54\) 个，没有涂漆面的有 \((5-2)^3=27\) 个.
所以
\[
E(X)=\frac{3\cdot8+2\cdot36+1\cdot54+0\cdot27}{125}=\frac{150}{125}=\frac65
\]
故选 B.
\end{solution}

\begin{problem}
已知 \(a\) 为常数，函数 \(f(x) = x(\ln x - ax)\) 有两个极值点 \(x_1\)，\(x_2\)（\(x_1 < x_2\)），则（\quad）

\choices
    {\(f(x_1) > 0\)，\(f(x_2) > -\frac{1}{2}\)}
    {\(f(x_1) < 0\)，\(f(x_2) < -\frac{1}{2}\)}
    {\(f(x_1) > 0\)，\(f(x_2) < -\frac{1}{2}\)}
    {\(f(x_1) < 0\)，\(f(x_2) > -\frac{1}{2}\)}
\end{problem}

\begin{answer}
D
\end{answer}

\begin{solution}
函数定义域为 \(x>0\)，且
\(f'(x)=\ln x+1-2ax\).
要使 \(f'(x)=0\) 有两个根，必须有 \(a>0\). 函数 \(h(x)=\ln x+1-2ax\) 在 \(x=\dfrac{1}{2a}\) 处取得最大值 \(-\ln(2a)\)，故两个根的存在要求 \(-\ln(2a)>0\)，即 \(0<a<\dfrac12\).
又 \(h(1)=1-2a>0\)，而 \(h(x)\) 在 \(x\to0^+\) 和 \(x\to+\infty\) 时均趋于负无穷，因此 \(x_1<1<x_2\).
在极值点处由 \(\ln x+1=2ax\) 得 \(ax=\dfrac{\ln x+1}{2}\)，于是
\(f(x)=x(\ln x-ax)=\dfrac{x}{2}(\ln x-1)\).
所以 \(f(x_1)<0\). 对 \(x>1\)，函数 \(x(\ln x-1)\) 的导数为 \(\ln x>0\)，且在 \(x=1\) 时其值为 \(-1\)，故 \(f(x_2)>-\dfrac12\).
因此选 D.
\end{solution}

\section{填空题}

\begin{problem}
从某小区抽取 100 户居民进行月用电量调查，发现其用电量都在 50 至 350 度之间，频率分布直方图如图所示.

\begin{enumerate}
\item 直方图中 \(x\) 的值为\fillinblank{}；
\item 在这些用户中，用电量落在区间 \([100,250)\) 内的户数为\fillinblank{}.
\end{enumerate}

\bitmapfigure[max width=0.45\linewidth]{img/2013/hubei_science/q11_fig1.png}
\end{problem}

\begin{answer}
\(0.0044\)；\(70\)
\end{answer}

\begin{solution}
每个组距为 \(50\)，频率分布直方图的各柱面积之和为 \(1\)，所以
\(50\left(0.0024+0.0036+0.0060+x+0.0024+0.0012\right)=1\).
解得 \(x=0.0044\).
用电量落在 \([100,250)\) 内的频率为
\(50(0.0036+0.0060+0.0044)=0.70\)，因此相应户数为 \(100\times0.70=70\) 户.
\end{solution}

\begin{problem}
阅读如图所示的程序框图，运行相应的程序，输出的结果 \(i=\fillinblank{}\).

\texfigure[max width=0.46\linewidth]{img_repaint/2013/hubei_science/q12_fig1.tex}
\end{problem}

\begin{answer}
5
\end{answer}

\begin{solution}
程序开始时 \(a=10\)，\(i=1\). 逐次执行循环：当 \(a=10\) 时，\(a\) 不是 \(4\) 且为偶数，令 \(a=10/2=5\)，再令 \(i=2\)；当 \(a=5\) 时，令 \(a=3\times5+1=16\)，再令 \(i=3\)；当 \(a=16\) 时，令 \(a=16/2=8\)，再令 \(i=4\)；当 \(a=8\) 时，令 \(a=8/2=4\)，再令 \(i=5\). 此时满足 \(a=4\)，输出 \(i=5\).
\end{solution}

\begin{problem}
设 \(x\)，\(y\)，\(z \in \R\)，且满足：\(x^2 + y^2 + z^2 = 1\)，\(x + 2y + 3z = \sqrt{14}\)，则 \(x + y + z =\fillinblank{}\).
\end{problem}

\begin{answer}
\(\frac{3\sqrt{14}}{7}\)
\end{answer}

\begin{solution}
由柯西不等式，
\((x+2y+3z)^2\leqslant(1^2+2^2+3^2)(x^2+y^2+z^2)=14\).
题设给出 \((x+2y+3z)^2=14\)，故等号成立. 于是 \((x,y,z)\) 与 \((1,2,3)\) 同向，设 \((x,y,z)=k(1,2,3)\). 代入 \(x+2y+3z=\sqrt{14}\)，得 \(14k=\sqrt{14}\)，所以 \(k=\dfrac{\sqrt{14}}{14}\).
因此 \(x+y+z=6k=\dfrac{3\sqrt{14}}{7}\).
\end{solution}

\begin{problem}
古希腊毕达哥拉斯学派的数学家研究过各种多边形数，如三角形数 \(1\)，\(3\)，\(6\)，\(10\)，\(\cdots\)，第 \(n\) 个三角形数为 \(\frac{n(n + 1)}{2} = \frac{1}{2}n^2 + \frac{1}{2}n\). 记第 \(n\) 个 \(k\) 边形数为 \(N(n, k)\)（\(k \geqslant 3\)），以下列出了部分 \(k\) 边形数中第 \(n\) 个数的表达式：

三角形数 \(N(n, 3) = \frac{1}{2}n^2 + \frac{1}{2}n\)，
正方形数 \(N(n, 4) = n^2\)，
五边形数 \(N(n, 5) = \frac{3}{2}n^2 - \frac{1}{2}n\)，
六边形数 \(N(n, 6) = 2n^2 - n\)，
……

可以推测 \(N(n, k)\) 的表达式，由此计算 \(N(10, 24) =\fillinblank{}\).
\end{problem}

\begin{answer}
1000
\end{answer}

\begin{solution}
观察题中各式，\(n^2\) 的系数在 \(k=3\)，\(4\)，\(5\)，\(6\) 时依次为 \(\dfrac12\)，\(1\)，\(\dfrac32\)，\(2\)，构成公差为 \(\dfrac12\) 的等差数列，故该系数为 \(\dfrac{k-2}{2}\). \(n\) 的系数依次为 \(\dfrac12\)，\(0\)，\(-\dfrac12\)，\(-1\)，故该系数为 \(\dfrac{4-k}{2}\). 因此
\(N(n,k)=\dfrac{k-2}{2}n^2+\dfrac{4-k}{2}n\).
代入 \(n=10\)、\(k=24\)，得
\(N(10,24)=\dfrac{22}{2}\cdot10^2+\dfrac{-20}{2}\cdot10=1100-100=1000\).
\end{solution}

\section{选考题：填空题}

\begin{problem}
如图，圆 \(O\) 上一点 \(C\) 在直径 \(AB\) 上的射影为 \(D\)，点 \(D\) 在半径 \(OC\) 上的射影为 \(E\)，若 \(AB = 3AD\)，则 \(\dfrac{CE}{EO}\) 的值为\fillinblank{}.

\bitmapfigure[max width=0.34\linewidth]{img/2013/hubei_science/q15_fig1.png}
\end{problem}

\begin{answer}
8
\end{answer}

\begin{solution}
设 \(AD=x\)，则 \(AB=3x\)，圆的半径为 \(OA=\dfrac{3x}{2}\). 因为 \(D\) 在线段 \(AB\) 上，所以
\(OD=OA-AD=\dfrac{x}{2}\).
在直角三角形 \(ODC\) 中，\(OC\) 是斜边，且 \(DE\perp OC\). 由射影定理，
\(OE\cdot OC=OD^2\)，从而 \(OE=\dfrac{(x/2)^2}{3x/2}=\dfrac{x}{6}\).
于是 \(CE=OC-OE=\dfrac{3x}{2}-\dfrac{x}{6}=\dfrac{4x}{3}\)，故
\(\dfrac{CE}{EO}=\dfrac{4x/3}{x/6}=8\).
\end{solution}

\begin{problem}
在直角坐标系 \(xOy\) 中，椭圆 \(C\) 的参数方程为 \(\begin{cases}
x = a\cos\varphi,\\
y = b\sin\varphi,
\end{cases}\)
（\(\varphi\) 为参数，\(a > b > 0\)），在极坐标系（与直角坐标系 \(xOy\) 取相同的长度单位，且以原点 \(O\) 为极点，以 \(x\) 轴正半轴为极轴）中，直线 \(l\) 与圆 \(O\) 的极坐标方程分别为 \(\rho\sin\left(\theta + \frac{\pi}{4}\right) = \frac{\sqrt{2}}{2}m\)（\(m\) 为非零常数）与 \(\rho = b\)，若直线 \(l\) 经过椭圆 \(C\) 的焦点，且与圆 \(O\) 相切，则椭圆 \(C\) 的离心率为\fillinblank{}.
\end{problem}

\begin{answer}
\(\frac{\sqrt{6}}{3}\)
\end{answer}

\begin{solution}
将直线的极坐标方程化为直角坐标方程. 由 \(x=\rho\cos\theta\)、\(y=\rho\sin\theta\)，有
\(\rho\sin\left(\theta+\dfrac{\pi}{4}\right)=\dfrac{x+y}{\sqrt{2}}\)，所以直线 \(l\) 的方程为 \(x+y=m\).
直线 \(l\) 与圆 \(O: x^2+y^2=b^2\) 相切，故原点到直线的距离满足 \(\dfrac{|m|}{\sqrt{2}}=b\)，即 \(|m|=\sqrt{2}b\).
椭圆焦点为 \((\pm c,0)\)，其中 \(c^2=a^2-b^2\). 由于直线 \(l\) 经过一个焦点，有 \(|m|=c\)，故 \(c=\sqrt{2}b\). 于是 \(a^2=b^2+c^2=3b^2\)，椭圆离心率为
\(e=\dfrac{c}{a}=\dfrac{\sqrt{2}b}{\sqrt{3}b}=\dfrac{\sqrt{6}}{3}\).
\end{solution}

\section{解答题}

\begin{problem}
在 \(\triangle ABC\) 中，角 \(A\)，\(B\)，\(C\) 对应的边分别是 \(a\)，\(b\)，\(c\). 已知 \(\cos 2A - 3\cos (B + C) = 1\).
\begin{enumerate}
\item 求角 \(A\) 的大小；
\item 若 \(\triangle ABC\) 的面积 \(S = 5\sqrt{3}\)，\(b = 5\)，求 \(\sin B \sin C\) 的值.
\end{enumerate}
\end{problem}

\begin{solution}
\begin{enumerate}
\item 因为 \(A+B+C=\pi\)，所以 \(B+C=\pi-A\)，从而 \(\cos(B+C)=-\cos A\). 代入已知等式，得 \(\cos 2A+3\cos A=1\).
令 \(u=\cos A\)，利用 \(\cos 2A=2\cos^2A-1\)，可得 \(2u^2-1+3u=1\)，即 \(2u^2+3u-2=0\)，从而 \((2u-1)(u+2)=0\).
因为 \(0<A<\pi\)，所以 \(-1<u<1\)，故 \(u=\frac{1}{2}\)，即 \(A=\frac{\pi}{3}\).

\item 由三角形面积公式 \(S=\frac{1}{2}bc\sin A\)，得 \(5\sqrt{3}=\frac{1}{2}\cdot5\cdot c\cdot\frac{\sqrt{3}}{2}\)，所以 \(c=4\). 由余弦定理，\(a^2=b^2+c^2-2bc\cos A=25+16-2\cdot5\cdot4\cdot\frac{1}{2}=21\). 由正弦定理，\(\sin B=\frac{b}{a}\sin A\)，\(\sin C=\frac{c}{a}\sin A\)，因此
\[
\sin B\sin C
=\frac{bc}{a^2}\sin^2A
=\frac{5\cdot4}{21}\cdot\frac{3}{4}
=\frac{5}{7}
\]
\end{enumerate}
\end{solution}

\begin{problem}
已知等比数列 \(\{a_{n}\}\) 满足：\(\left|a_2 - a_3\right| = 10\)，\(a_1a_2a_3 = 125\).
\begin{enumerate}
\item 求数列 \(\{a_{n}\}\) 的通项公式；
\item 是否存在正整数 \(m\)，使得 \(\frac{1}{a_1} + \frac{1}{a_2} + \cdots + \frac{1}{a_m} \geqslant 1\). 若存在，求 \(m\) 的最小值；若不存在，说明理由.
\end{enumerate}
\end{problem}

\begin{solution}
设等比数列的公比为 \(q\). 由 \(a_2=a_1q\)，有
\[
a_1a_2a_3=a_1\cdot a_1q\cdot a_1q^2=(a_1q)^3=a_2^3
\]
因此 \(a_2^3=125\)，从而 \(a_2=5\). 又 \(|a_2-a_3|=|5-5q|=10\)，即 \(|1-q|=2\)，所以 \(q=3\) 或 \(q=-1\).

当 \(q=3\) 时，\(a_1=\frac{a_2}{q}=\frac{5}{3}\)，故
\(a_n=\frac{5}{3}\cdot3^{n-1}=5\cdot3^{n-2}\).
当 \(q=-1\) 时，\(a_1=-5\)，故
\(a_n=-5(-1)^{n-1}=5(-1)^n\). 因此数列的通项公式有两种可能：\(a_n=5\cdot3^{n-2}\) 或 \(a_n=5(-1)^n\).

对第（2）问分别讨论.

当 \(q=3\) 时，
\[
\sum_{k=1}^{m}\frac{1}{a_k}
=\frac{3}{5}\sum_{k=0}^{m-1}\left(\frac{1}{3}\right)^k
=\frac{9}{10}\left(1-\frac{1}{3^m}\right)<\frac{9}{10}<1
\]
当 \(q=-1\) 时，\(a_n=5(-1)^n\)，所以
\[
\sum_{k=1}^{m}\frac{1}{a_k}
=\frac{1}{5}\sum_{k=1}^{m}(-1)^k
=\begin{cases}
0,&m\text{ 为偶数},\\
-\frac{1}{5},&m\text{ 为奇数}.
\end{cases}
\]
这两种情形下的和都不可能不小于 1，故不存在满足条件的正整数 \(m\).
\end{solution}

\begin{problem}
如图，\(AB\) 是圆 \(O\) 的直径，点 \(C\) 是圆 \(O\) 上异于 \(A\)，\(B\) 的点，直线 \(PC \perp\) 平面 \(ABC\)，\(E\)，\(F\) 分别是 \(PA\)，\(PC\) 的中点.
\begin{enumerate}
\item 记平面 \(BEF\) 与平面 \(ABC\) 的交线为 \(l\)，试判断直线 \(l\) 与平面 PAC 的位置关系，并加以证明；
\item 设 （1） 中的直线 \(l\) 与圆 \(O\) 的另一个交点为 \(D\)，且点 \(Q\) 满足 \(\overrightarrow{DQ} = \frac{1}{2}\overrightarrow{CP}\)，记直线 \(PQ\) 与平面 \(ABC\) 所成的角为 \(\theta\)，异面直线 \(PQ\) 与 \(EF\) 所成的角为 \(\alpha\)，二面角 \(E - l - C\) 的大小为 \(\beta\)，求证： \(\sin \theta = \sin \alpha \sin \beta\).
\end{enumerate}
\end{problem}

\begin{solution}
\begin{enumerate}
\item 因为 \(AB\) 是圆的直径，所以 \(\angle ACB=90^{\circ}\)，即 \(AC\perp BC\). 建立空间直角坐标系，使 \(C\) 为原点，\(AC\)、\(BC\)、\(PC\) 分别为 \(x\) 轴、\(y\) 轴、\(z\) 轴，设 \(A=(a,0,0)\)，\(B=(0,b,0)\)，\(P=(0,0,h)\)，其中 \(a>0\)，\(b>0\)，\(h>0\).
则 \(E=\left(\frac{a}{2},0,\frac{h}{2}\right)\)，\(F=\left(0,0,\frac{h}{2}\right)\).
由于 \(EF\parallel AC\)，平面 \(BEF\) 与平面 \(ABC\) 的交线 \(l\) 经过公共点 \(B\)，且平行于 \(AC\).

平面 \(PAC\) 与平面 \(ABC\) 的交线为 \(AC\)，而 \(B\notin AC\)，所以直线 \(l\) 不与平面 \(PAC\) 相交，又因 \(l\parallel AC\)，可得 \(l\parallel\text{平面 }PAC\).

\item 在上述坐标系中，圆 \(O\) 经过 \(A\)、\(B\)、\(C\)，其方程为
\(x^2+y^2-ax-by=0\).
由第（1）问，直线 \(l\) 的方程为 \(y=b\)，\(z=0\). 将其代入圆的方程，得 \(x^2-ax=0\)，所以除点 \(B=(0,b,0)\) 外，另一个交点为 \(D=(a,b,0)\). 由 \(\overrightarrow{DQ}=\frac{1}{2}\overrightarrow{CP}\)，得 \(Q=\left(a,b,\frac{h}{2}\right)\).
于是可取 \(\overrightarrow{PQ}=\left(a,b,-\frac{h}{2}\right)\)，\(\overrightarrow{EF}=\left(-\frac{a}{2},0,0\right)\). 令 \(L=\sqrt{a^2+b^2+\frac{h^2}{4}}\). 直线 \(PQ\) 与平面 \(ABC\) 所成的角为 \(\theta\)，故 \(\sin\theta=\frac{h/2}{L}\). 直线 \(EF\) 的方向为 \(x\) 轴方向，因此 \(\cos\alpha=\frac{a}{L}\)，\(\sin\alpha=\frac{\sqrt{b^2+h^2/4}}{L}\).

二面角 \(E-l-C\) 的一个面为平面 \(E-l\)，另一个面为平面 \(C-l=ABC\). 在垂直于 \(l\) 的截面内，平面 \(E-l\) 的方向向量可取 \((-b,h/2)\)，所以 \(\sin\beta=\frac{h/2}{\sqrt{b^2+h^2/4}}\).
从而
\[
\sin\alpha\sin\beta
=\frac{\sqrt{b^2+h^2/4}}{L}\cdot
\frac{h/2}{\sqrt{b^2+h^2/4}}
=\frac{h/2}{L}
=\sin\theta
\]
故 \(\sin\theta=\sin\alpha\sin\beta\).
\end{enumerate}
\end{solution}

\begin{problem}
假设每天从甲地去乙地的旅客人数 \(X\) 是服从正态分布 \(N(800,50^{2})\) 的随机变量，记一天中从甲地去乙地的旅客人数不超过 900 的概率为 \(P_{0}\).
\begin{enumerate}
\item 求 \(P_{0}\) 的值；（参考数据：若 \(X \sim N(\mu, \sigma^{2})\)，有 \(P(\mu - \sigma < X \leqslant \mu + \sigma) = 0.6826\)，\(P(\mu - 2\sigma < X \leqslant \mu + 2\sigma) = 0.9544\)，\(P(\mu - 3\sigma < X \leqslant \mu + 3\sigma) = 0.9974\)）
\item 某客运公司用 \(A\)，\(B\) 两种型号的车辆承担甲，乙两地间的长途客运业务，每车每天往返一次，\(A\)，\(B\) 两种车辆的载客量分别为 36 人和 60 人，从甲地去乙地的营运成本分别为 1600 元/辆和 2400 元/辆. 公司拟组建一个不超过 21 辆车的客运车队，并要求 \(B\) 型车不多于 \(A\) 型车 7 辆. 若每天要以不小于 \(p_0\) 的概率运完从甲地去乙地的旅客，且使公司从甲地去乙地的营运成本最小，那么应配备 \(A\) 型车，\(B\) 型车各多少辆？
\end{enumerate}
\end{problem}

\begin{solution}
\begin{enumerate}
\item 因为 \(900=800+2\times50\)，所以 \(P_0=P(X\leqslant800+2\times50)\).
正态分布关于均值对称，且 \(P(700<X\leqslant900)=0.9544\).
因此
\[
P_0=P(X\leqslant800)+P(800<X\leqslant900)
=0.5+\frac{0.9544}{2}=0.9772
\]

\(p_0\) 按题意取为第一问所得的 \(P_0\)，即 \(p_0=0.9772\).

\item 设配备 \(x\) 辆 \(A\) 型车、\(y\) 辆 \(B\) 型车，其中 \(x\)，\(y\) 为非负整数. 要使运力满足概率要求，车辆总载客量至少为 900 人，因此 \(36x+60y\geqslant900\)，即 \(3x+5y\geqslant75\). 另外，\(x+y\leqslant21\)，\(y\leqslant x+7\). 营运成本为 \(1600x+2400y=800(2x+3y)\).

下面证明 \(2x+3y\) 不小于 46. 若 \(2x+3y\leqslant45\)，则由运力约束
\[
75\leqslant3x+5y
=\frac{3}{2}(2x+3y)+\frac{1}{2}y
\leqslant67.5+\frac{1}{2}y
\]
可得 \(y\geqslant15\). 又因 \(y\leqslant x+7\)，所以 \(x\geqslant8\)，从而 \(2x+3y\geqslant2\cdot8+3\cdot15=61\)，这与 \(2x+3y\leqslant45\) 矛盾. 因此 \(2x+3y\geqslant46\).

取 \(x=5\)，\(y=12\)，则 \(36x+60y=36\cdot5+60\cdot12=900\)，\(x+y=17\leqslant21\)，\(y=x+7\)，并且 \(1600x+2400y=1600\cdot5+2400\cdot12=36800\).
所以应配备 \(A\) 型车 5 辆，\(B\) 型车 12 辆，最小营运成本为 36800 元.
\end{enumerate}
\end{solution}

\begin{problem}
如图，已知椭圆 \(C_1\) 与 \(C_2\) 的中心在坐标原点 \(O\)，长轴均为 \(MN\) 且在 \(x\) 轴上，短轴长分别为 \(2m\)，\(2n\)（\(m > n\)），过原点且不与 \(x\) 轴重合的直线 \(l\) 与 \(C_1\)，\(C_2\) 的四个交点按纵坐标从大到小依次为 \(A\)，\(B\)，\(C\)，\(D\)，记 \(\lambda = \frac{m}{n}\)，\(\triangle BDM\) 和 \(\triangle ABN\) 的面积分别为 \(S_1\) 和 \(S_2\).
\begin{enumerate}
\item 当直线 \(l\) 与 \(y\) 轴重合时，若 \(S_{1} = \lambda S_{2}\)，求 \(\lambda\) 的值；
\item 当 \(\lambda\) 变化时，是否存在与坐标轴不重合的直线 \(l\) 使得 \(S_{1} = \lambda S_{2}\) 并说明理由.
\end{enumerate}
\end{problem}

\begin{solution}
设 \(OM=ON=a\)，则 \(a>m>n>0\).

\begin{enumerate}
\item 当 \(l\) 与 \(y\) 轴重合时，
有 \(A=(0,m)\)，\(B=(0,n)\)，\(C=(0,-n)\)，\(D=(0,-m)\)，\(M=(-a,0)\)，\(N=(a,0)\). 因此 \(S_1=\frac{1}{2}a(m+n)\)，\(S_2=\frac{1}{2}a(m-n)\). 由 \(S_1=\lambda S_2\) 及 \(\lambda=\frac{m}{n}\)，得 \(m+n=\frac{m}{n}(m-n)\).
令 \(m=\lambda n\)，可得 \(\lambda+1=\lambda(\lambda-1)\)，即 \(\lambda^2-2\lambda-1=0\).
因为 \(\lambda>1\)，所以 \(\lambda=1+\sqrt{2}\).

\item 设直线 \(l\) 与正半轴方向的夹角为 \(\varphi\)，取直线在上半平面的射线为正方向. 记该射线与 \(C_1\)，\(C_2\) 的交点到原点的距离分别为 \(R_m\)，\(R_n\)，则
\[
R_m=\left(\frac{\cos^2\varphi}{a^2}+\frac{\sin^2\varphi}{m^2}\right)^{-1/2}
\]
并且
\[
R_n=\left(\frac{\cos^2\varphi}{a^2}+\frac{\sin^2\varphi}{n^2}\right)^{-1/2}
\]
记上半平面内的两个交点为 \(A=(x_m,y_m)\)、\(B=(x_n,y_n)\)，则
有 \(y_m=R_m\sin\varphi\)，\(y_n=R_n\sin\varphi\)，\(C=-B\)，\(D=-A\).
由三角形面积的行列式公式，
\[
S_1=\frac{1}{2}\left|\overrightarrow{BD}\times\overrightarrow{BM}\right|
=\frac{1}{2}a(y_m+y_n)
\]
以及
\[
S_2=\frac{1}{2}\left|\overrightarrow{AB}\times\overrightarrow{AN}\right|
=\frac{1}{2}a(y_m-y_n)
\]
令 \(\rho=\frac{y_m}{y_n}=\frac{R_m}{R_n}\)，则条件 \(S_1=\lambda S_2\) 等价于 \(\frac{\rho+1}{\rho-1}=\lambda\)，即 \(\rho=\frac{\lambda+1}{\lambda-1}\).

当直线不与坐标轴重合时，\(\sin\varphi\neq0\) 且 \(\cos\varphi\neq0\). 由
\[
\rho^2=\lambda^2\frac{n^2\cos^2\varphi+a^2\sin^2\varphi}{m^2\cos^2\varphi+a^2\sin^2\varphi}
\]
可知 \(1<\rho<\lambda\). 因此必要条件为 \(\frac{\lambda+1}{\lambda-1}<\lambda\)，即 \(\lambda^2-2\lambda-1>0\).
结合 \(\lambda>1\)，得到 \(\lambda>1+\sqrt{2}\).

反之，当 \(\lambda>1+\sqrt{2}\) 时，令 \(m=\lambda n\)，由上式令 \(\rho=\frac{\lambda+1}{\lambda-1}\)，整理得
\[
\tan^2\varphi
=\frac{4\lambda^3n^2}{a^2(\lambda^2+1)(\lambda^2-2\lambda-1)}
\]
右端为正且有限，故存在既不与 \(x\) 轴也不与 \(y\) 轴重合的直线 \(l\) 满足条件. 综上，当且仅当 \(\lambda>1+\sqrt{2}\) 时，存在满足条件的直线 \(l\).
\end{enumerate}
\end{solution}

\begin{problem}
设 \(n\) 为正整数，\(r\) 为正有理数.
\begin{enumerate}
\item 求函数 \( f(x) = (1 + x)^{r + 1} - (r + 1)x - 1 \)（\(x > -1\)）的最小值；
\item 证明： \(\frac{n^{r+1} - (n-1)^{r+1}}{r+1} < n^{r} < \frac{(n+1)^{r+1} - n^{r+1}}{r+1}\)；
\item 设 \(x \in \R\)，记 \([x]\) 为不小于 \(x\) 的最小整数，例如 \([2] = 2\)，\([\pi] = 4\)，\(\left[\frac{3}{-2}\right] = -1\). 令 \(S = \sqrt[3]{81} + \sqrt[3]{82} + \sqrt[3]{83} + \cdots + \sqrt[3]{125}\)，求 \([S]\) 的值.（参考数据： \(80^{\frac{4}{3}} \approx 344.7\)，\(81^{\frac{4}{3}} \approx 350.5\)，\(124^{\frac{4}{3}} \approx 618.3\)，\(126^{\frac{4}{3}} \approx 631.7\)）
\end{enumerate}
\end{problem}

\begin{solution}
\begin{enumerate}
\item 令 \(t=1+x\)，则 \(t>0\)，且 \(f(x)=t^{r+1}-(r+1)t+r\). 于是 \(f'(x)=(r+1)(t^r-1)=(r+1)((1+x)^r-1)\).
当 \(-1<x<0\) 时，\(0<1+x<1\)，所以 \((1+x)^r<1\)，从而 \(f'(x)<0\)；当 \(x>0\) 时，\((1+x)^r>1\)，从而 \(f'(x)>0\). 因此函数在 \(x=0\) 处取得最小值，且
\(f(0)=1-0-1=0\).

\item 当 \(n=1\) 时，\(\frac{1^{r+1}-0^{r+1}}{r+1}=\frac{1}{r+1}<1=1^r\).

当 \(n\geqslant2\) 时，将 \(x=-\frac{1}{n}\) 代入第（1）问的结论. 由于 \(-\frac{1}{n}\neq0\)，有 \(\left(1-\frac{1}{n}\right)^{r+1}+\frac{r+1}{n}-1>0\). 两边同乘以正数 \(n^{r+1}\)，得 \((n-1)^{r+1}+(r+1)n^r-n^{r+1}>0\)，即 \(\frac{n^{r+1}-(n-1)^{r+1}}{r+1}<n^r\).

对任意正整数 \(n\)，将 \(x=\frac{1}{n}\) 代入，同样因为 \(x\neq0\)，有 \(\left(1+\frac{1}{n}\right)^{r+1}-\frac{r+1}{n}-1>0\). 两边同乘以正数 \(n^{r+1}\)，得 \((n+1)^{r+1}-(r+1)n^r-n^{r+1}>0\)，即 \(n^r<\frac{(n+1)^{r+1}-n^{r+1}}{r+1}\).
两式合并即得所证不等式.

\item 在第（2）问中取 \(r=\frac{1}{3}\)，对 \(n=81\)，\(82\)，\(\ldots\)，\(125\) 分别求和，得到
\[
\frac{3}{4}\left(125^{4/3}-80^{4/3}\right)
<S<
\frac{3}{4}\left(126^{4/3}-81^{4/3}\right)
\]
因为 \(125^{4/3}=625\)，结合题中参考数据，\(\frac{3}{4}(625-344.7)\approx210.2<S\)，且 \(S<\frac{3}{4}(631.7-350.5)\approx210.9\).
所以 \(210<S<211\). 题目中的 \([S]\) 表示不小于 \(S\) 的最小整数，故 \([S]=211\).
\end{enumerate}
\end{solution}
